\chapter{Playlists}\label{chap:listes}

\section{Multiple playlists}
\soft handles any number of playlists, plus two special ones:
\begin{description}
  \item[Default playlist] created at first launch, it takes over your previous playlist.
  \item[Favourites] filled by the \textbf{$\equiv$+} button $\rightarrow$ ``Add to favorites''.
\end{description}
The displayed playlist is also the played one. Switching in the selector saves the playlist left and loads the new
one; the current track keeps playing and ``next'' follows the new playlist.

\capture{playlist-en-top}{Playlist selector}

\section{Managing playlists}
The selector menu offers: switch, new playlist, rename, duplicate, delete, export as M3U, import M3U. Deleting the
displayed playlist switches to another one.

\section{Adding tracks}
\begin{itemize}
  \item From Now playing: \textbf{$\equiv$+} $\rightarrow$ ``Add to playlist\dots'', then a playlist or ``New
        playlist\dots''.
  \item From the browser: selection $\rightarrow$ Add.
  \item From a playlist: selection $\rightarrow$ ``Add to playlist\dots'' copies tracks to another playlist.
\end{itemize}

\section{Play next queue}\label{sec:file}
Queued tracks play after the current one, in the order they were added. Then playback resumes where it was in the
playlist or folder. Choosing another track yourself keeps the queue but drops that resume point.

\section{M3U import and export}
Export writes the displayed playlist as extended M3U (title, author, duration). Locations are complete, including
inside archives and online catalogues, so an exported playlist is re-imported identically. Import also accepts M3U
from other players (plain paths) and creates a new playlist named after the file. Former XSPF playlists remain
available in the browser.

\section{Sort, move, statistics}
Sorting is by title, author or duration. Each track's handle moves it. Statistics give track count, distinct files
and total duration of the playlist or the selection.
